%% ============================================================
%%  OLIBOT — Guión del ponente (presentación dirección escolar)
%%  Compilar: pdflatex guion_presentacion_colegio.tex
%% ============================================================
\documentclass[11pt, a4paper]{article}

\usepackage[spanish]{babel}
\usepackage[utf8]{inputenc}
\usepackage[T1]{fontenc}
\usepackage{lmodern}
\usepackage{geometry}
\usepackage{xcolor}
\usepackage{enumitem}
\usepackage{titlesec}
\usepackage{booktabs}
\usepackage{mdframed}
\usepackage{microtype}
\usepackage{parskip}
\usepackage{fontawesome5}

\geometry{left=2.2cm, right=2.2cm, top=2.5cm, bottom=2.5cm}

\definecolor{verde}   {RGB}{34,  139,  86}
\definecolor{azul}    {RGB}{30,   90, 160}
\definecolor{naranja} {RGB}{230, 120,  30}
\definecolor{gris}    {RGB}{100, 100, 100}

%% Entornos
\newmdenv[
  linecolor=verde, linewidth=1.5pt,
  backgroundcolor=verde!6,
  innerleftmargin=10pt, innerrightmargin=10pt,
  innertopmargin=6pt, innerbottommargin=6pt,
  skipabove=6pt, skipbelow=4pt
]{decir}

\newmdenv[
  linecolor=naranja, linewidth=1pt,
  topline=false, bottomline=false, rightline=false,
  backgroundcolor=naranja!5,
  innerleftmargin=10pt, innerrightmargin=10pt,
  innertopmargin=5pt, innerbottommargin=5pt,
  skipabove=4pt, skipbelow=4pt
]{nota}

\newcommand{\diapositiva}[2]{%
  \vspace{8pt}
  \noindent{\color{azul}\rule{\linewidth}{1.5pt}}\\
  \noindent{\large\bfseries\color{azul} Diapositiva #1 — #2}
  \vspace{4pt}\\
}

\newcommand{\tiempo}[1]{%
  {\footnotesize\color{gris}\faIcon{clock}\; \textit{Tiempo estimado: #1}}
  \smallskip\\
}

\title{%
  {\large OLIBOT — Propuesta de evaluación}\\[4pt]
  {\LARGE\bfseries Guión del Ponente}\\[4pt]
  {\normalsize Para presentación a la dirección del centro escolar}
}
\author{María Emilia Núñez Guerrero}
\date{Junio 2026}

\begin{document}
\maketitle
\tableofcontents
\newpage

%% ============================================================
\section*{Antes de entrar}
\addcontentsline{toc}{section}{Antes de entrar}
%% ============================================================

\begin{nota}
\textbf{Preparación previa (5 min antes):}
\begin{itemize}[nosep]
  \item Abrir la tablet con OLIBOT cargado y listo en la pantalla de selección de estudiante.
  \item Tener el proyector conectado con las diapositivas en la primera.
  \item Cerrar cualquier otra pestaña del navegador.
  \item Asegurarse de que el audio del dispositivo funciona correctamente para la demo.
  \item Respirar, sonreír. El objetivo es crear una conversación, no vender un producto.
\end{itemize}
\end{nota}

%% ============================================================
\section{Bienvenida e introducción (Diap. 1--2)}
%% ============================================================

\diapositiva{1}{Portada}
\tiempo{1 minuto}

\begin{decir}
«Buenos días / tardes. Mi nombre es María Emilia Núñez Guerrero, soy estudiante del Máster
en Inteligencia Artificial de la Universidad Carlos III de Madrid.

Hoy vengo a presentarles OLIBOT, un asistente educativo digital que estoy desarrollando
como Trabajo de Fin de Máster, y que me gustaría tener la oportunidad de evaluar
con los alumnos de su centro.

La presentación es breve — unos veinte minutos — y al final haremos una demostración
en directo para que puedan ver exactamente con qué interactuarían los niños.»
\end{decir}

\diapositiva{2}{Índice}
\tiempo{30 segundos}

\begin{decir}
«Les voy a contar cuatro cosas: el porqué de este proyecto, qué es y qué hace OLIBOT,
cómo planteo la evaluación, y qué garantías ofrezco al centro. Y para cerrar, la demostración.»
\end{decir}

%% ============================================================
\section{El contexto educativo (Diap. 3--4)}
%% ============================================================

\diapositiva{3}{El currículo de Infantil en la Comunidad de Madrid}
\tiempo{2 minutos}

\begin{decir}
«El currículo de Educación Infantil de la Comunidad de Madrid, el Decreto 36/2022,
establece que la iniciación a la lectoescritura es un eje central desde los tres años:
grafomotricidad, conciencia fonológica, reconocimiento de letras y números.

Y, al mismo tiempo, todos sabemos que en el aula de Infantil hay hasta 25 niños con
ritmos de aprendizaje muy distintos. Atender individualmente a cada uno mientras
se lleva una clase completa es, sencillamente, muy difícil.

La pregunta que me hice fue: ¿puede la tecnología ayudar a los docentes a dar
esa atención más personalizada, respetando lo que marca el currículo?»
\end{decir}

\begin{nota}
Si la directora interrumpe con una pregunta aquí sobre el currículo, es una
\textbf{buena señal}: está comprometida. Responder con calma y volver al hilo.
\end{nota}

\diapositiva{4}{Los niños de hoy y la tecnología}
\tiempo{1 minuto 30 segundos}

\begin{decir}
«Los niños de tres a cinco años ya interactúan con tablets y asistentes de voz
en casa a diario. El reto no es si usar la tecnología — ya la usan — sino cómo usarla
con sentido educativo.

Lo que diferencia a OLIBOT de una aplicación de juegos cualquiera es que no tiene
publicidad, no recopila datos personales, y cada actividad está diseñada para alcanzar
un objetivo curricular concreto. No es entretenimiento: es práctica guiada.»
\end{decir}

%% ============================================================
\section{¿Qué es OLIBOT? (Diap. 5--7)}
%% ============================================================

\diapositiva{5}{OLIBOT en pocas palabras}
\tiempo{2 minutos}

\begin{decir}
«OLIBOT es un asistente que habla con el niño por voz. No hace falta saber leer
ni escribir para interactuar con él — que es precisamente el contexto de Infantil.

Al comenzar, el docente selecciona el perfil del niño: su nombre y su edad. A partir de ahí,
OLIBOT propone la actividad adecuada para ese nivel — vocales para los de tres años,
consonantes y sílabas para los de cinco — y va ajustando la dificultad según cómo le vaya.

Funciona en cualquier tablet con Chrome o Safari. No requiere instalación, no hay
cuenta de usuario, no hay contraseña.»
\end{decir}

\begin{nota}
Señalar el diagrama de la derecha: «Este ciclo es lo que experimenta el niño en cada
actividad: habla, OLIBOT elige, el niño traza o responde, OLIBOT celebra, avanza al
siguiente reto. Y vuelve a empezar.»
\end{nota}

\diapositiva{6}{¿Qué hace el niño durante la sesión?}
\tiempo{1 minuto 30 segundos}

\begin{decir}
«Hay tres tipos de actividades. La principal es el trazado de letras en pantalla táctil:
OLIBOT primero demuestra cómo se traza la letra — el niño ve un cursor rojo que recorre
el trazo — y luego es el niño quien lo intenta con el dedo.

Toda la interacción es por voz: OLIBOT habla, el niño escucha y responde hablando o
trazando. Y de vez en cuando, como descanso, hay fichas de colorear digitales.

Las sesiones duran entre quince y veinte minutos. La maestra permanece en el aula
en todo momento.»
\end{decir}

\diapositiva{7}{Principios pedagógicos}
\tiempo{2 minutos}

\begin{decir}
«El diseño pedagógico se basa en el concepto de andamiaje de Vygotsky: cuando el niño
tiene dificultades, OLIBOT no le da la respuesta directamente. Primero le anima a
intentarlo de nuevo; si sigue sin lograrlo, le da una pista específica; y solo si es necesario,
le hace la demostración completa. Siempre en positivo, nunca con reproches.

Cada niño tiene su perfil guardado, así que si un día le cuesta la letra M, al día siguiente
OLIBOT la vuelve a proponer para reforzarla.

En cuanto a seguridad: el programa solo responde sobre las actividades.
No puede hablar de nada ajeno a la lección. No graba el audio del niño.
No accede a internet durante la sesión.»
\end{decir}

%% ============================================================
\section{Plan de evaluación (Diap. 8--9)}
%% ============================================================

\diapositiva{8}{¿Qué queremos evaluar?}
\tiempo{2 minutos}

\begin{decir}
«El objetivo de la evaluación es medir si OLIBOT realmente ayuda. Para eso he definido
seis indicadores:

Primero, el \textbf{tiempo de atención sostenida}: cuánto tiempo permanece el niño
activamente implicado con la herramienta.

Segundo, el \textbf{rendimiento comparado}: comparar cuántas actividades supera el
grupo que usa OLIBOT frente al grupo que hace las mismas fichas en papel.

Tercero, la \textbf{intervención del adulto}: contar las veces que la maestra o yo
necesitamos intervenir para ayudar al niño.

Cuarto, la \textbf{satisfacción subjetiva}: al terminar, le mostramos al niño una
cara feliz y una triste y nos dice cómo se ha sentido.

Quinto, la \textbf{velocidad de progresión}: cuántas letras o sílabas nuevas supera
cada niño por sesión.

Y sexto, el \textbf{engagement voluntario}: cuántas veces el niño pide repetir
una actividad sin que se lo pidamos.»
\end{decir}

\begin{nota}
\textbf{Punto clave a enfatizar:} todos estos datos se recogen de forma anónima.
No hay ningún dato que permita identificar individualmente a un alumno.
\end{nota}

\diapositiva{9}{Diseño de la evaluación}
\tiempo{1 minuto 30 segundos}

\begin{decir}
«Propongo entre tres y cinco sesiones de quince a veinte minutos con un grupo
reducido de cuatro a ocho alumnos de cuatro o cinco años, divididos en dos subgrupos:
uno usa OLIBOT y el otro hace la misma actividad en papel. Las sesiones se realizan
en el momento que la tutora considere más oportuno — rincones, trabajo por grupos,
o tiempo libre guiado. En ningún caso interrumpiremos la dinámica de clase.

Material necesario: una sola tablet y conexión WiFi. Si no hay WiFi disponible,
nosotros aportamos el punto de acceso.

Al finalizar, les entrego un informe completo en una semana.»
\end{decir}

%% ============================================================
\section{Garantías y compromisos (Diap. 10--11)}
%% ============================================================

\diapositiva{10}{Lo que garantizamos}
\tiempo{1 minuto 30 segundos}

\begin{decir}
«Quiero ser muy clara en lo que respecta a la privacidad. Los perfiles de los niños
usan apodos o iniciales — sin nombres completos, sin fotos. El reconocimiento de voz
se procesa en el propio dispositivo: no se envía audio a ningún servidor externo.

Este trabajo es exclusivamente académico — parte de mi TFM de la Universidad Carlos III.
No hay ninguna empresa detrás, no hay uso comercial de los resultados.

A cambio de su colaboración, el centro recibe:
un informe anónimo del progreso de cada alumno,
un informe para el equipo docente con sugerencias,
y la posibilidad de seguir usando OLIBOT de forma completamente gratuita.

Si lo desean, firmamos una carta de compromiso formal antes de comenzar.»
\end{decir}

\diapositiva{11}{Lo que necesitamos de vosotros}
\tiempo{1 minuto}

\begin{decir}
«En concreto, solo les pido cinco cosas:
tres a cinco sesiones en horario escolar,
la presencia de la tutora durante las sesiones,
acceso WiFi o aceptar que ponga yo un punto de acceso,
la autorización de la dirección y de las familias — les facilitaré un modelo estándar —
y un cuestionario breve al docente al terminar.

Nada más. No modifica la dinámica del aula. La tutora puede en cualquier momento
decirme que paremos.»
\end{decir}

%% ============================================================
\section{Cierre y demostración (Diap. 12--14)}
%% ============================================================

\diapositiva{12}{¿Por qué vale la pena probar OLIBOT?}
\tiempo{1 minuto}

\begin{decir}
«Para cerrar la parte teórica: ¿qué gana el aula?

Mientras OLIBOT acompaña a un grupo de dos o tres niños, la maestra tiene tiempo
para atender individualmente a quienes más lo necesitan.
Cada niño avanza a su ritmo, sin presión ni comparaciones.
Las actividades refuerzan exactamente lo que se trabaja en clase.
Y los datos del programa pueden servir de apoyo a la evaluación continua del docente.

La tecnología, bien usada, no reemplaza a la maestra: multiplica su impacto.»
\end{decir}

\begin{nota}
\textbf{Pausa.} Mirar a la directora. Si hay algún gesto de duda, preguntar:
«¿Tiene alguna pregunta antes de ver la demostración?»
\end{nota}

\diapositiva{13}{DEMO EN VIVO}
\tiempo{5 minutos}

\begin{decir}
«¿Les parece bien si lo vemos en directo? Les voy a mostrar exactamente lo que
vería un niño en su primera sesión con OLIBOT.»
\end{decir}

\begin{nota}
\textbf{Checklist de la demo (en este orden):}
\begin{enumerate}[nosep]
  \item Tomar la tablet. Abrir OLIBOT en el navegador.
  \item Crear un perfil de niño: por ejemplo «Oli», 4 años.
  \item Mostrar la bienvenida y cómo OLIBOT saluda por voz.
  \item Esperar a que aparezca la actividad de trazado (vocal A o E).
  \item \textbf{Mostrar la demo automática}: «Fijaos, primero OLIBOT hace
        la demostración del trazo con ese cursor rojo.»
  \item Trazar la letra con el dedo para mostrar la interacción táctil.
  \item Demostrar la interacción por voz: hablar con OLIBOT.
  \item Si hay tiempo: mostrar el nivel fácil (con guía completa)
        vs. nivel difícil (solo los puntos de inicio y fin).
\end{enumerate}

\textbf{Frases útiles durante la demo:}\\
«Aquí ven la demostración automática — el niño siempre ve primero cómo se hace.»\\
«Si el niño se equivoca varias veces, OLIBOT aumenta la guía.»\\
«El docente puede ver el progreso aquí, en este panel lateral.»
\end{nota}

\diapositiva{14}{Próximos pasos}
\tiempo{1 minuto}

\begin{decir}
«Si les parece bien, les propongo estos pasos:

Hoy mismo pueden comentarlo con el equipo docente.
Esta semana podemos firmar la carta de colaboración y preparar
las autorizaciones para las familias, que yo les facilito.
En las próximas semanas acordamos las fechas de sesiones.
Y tras las sesiones, una semana después, recibirán el informe completo.

¿Tienen alguna pregunta?»
\end{decir}

\begin{nota}
\textbf{Si dicen que necesitan consultarlo:}\\
«Por supuesto, no hay prisa. ¿Puedo dejarles mi contacto y un pequeño
resumen por escrito?» → dejar tarjeta o anotar email.

\textbf{Si preguntan por la privacidad:}\\
«Los datos de los niños se almacenan únicamente como número de intentos
y tiempo de sesión, sin audio y sin nombre. Les puedo mostrar exactamente
qué registra el sistema si lo desean.»

\textbf{Si preguntan si el programa «hace todo el trabajo del docente»:}\\
«No, en absoluto. OLIBOT no evalúa, no observa al niño en su globalidad,
no detecta dificultades emocionales. La tutora sigue siendo imprescindible.
OLIBOT es como un cuaderno de práctica interactivo, no un sustituto.»

\textbf{Si preguntan por el coste:}\\
«Para el centro es completamente gratuito. Este es un proyecto académico,
no comercial.»
\end{nota}

%% ============================================================
\section*{Resumen de tiempos}
\addcontentsline{toc}{section}{Resumen de tiempos}
%% ============================================================

\vspace{8pt}
\begin{center}
\begin{tabular}{lll}
  \toprule
  \textbf{Bloque} & \textbf{Diapositivas} & \textbf{Tiempo} \\
  \midrule
  Bienvenida e índice    & 1--2  & 1,5 min \\
  Contexto educativo     & 3--4  & 3,5 min \\
  Qué es OLIBOT          & 5--7  & 5,5 min \\
  Plan de evaluación     & 8--9  & 3,5 min \\
  Garantías y compromisos & 10--11 & 2,5 min \\
  Cierre                 & 12    & 1 min   \\
  Demo en vivo           & 13    & 5 min   \\
  Próximos pasos y dudas & 14    & 2,5 min \\
  \midrule
  \textbf{Total} & & \textbf{$\approx$ 25 min} \\
  \bottomrule
\end{tabular}
\end{center}

\vspace{12pt}
\begin{nota}
\textbf{Consejo general:}\\
El tono debe ser cercano y honesto. Si algo no funciona en la demo, decirlo con
naturalidad: «La tecnología tiene estos momentos — así que en las sesiones reales
siempre estoy presente para resolver cualquier incidencia.»\\[4pt]
El objetivo no es impresionar con tecnología: es \textbf{transmitir confianza}.
\end{nota}

\end{document}
